\documentclass[11pt]{article}
% Shared preamble for the printable mini lessons (1.4 onward).
%   \documentclass[11pt]{article}
%   % Shared preamble for the printable mini lessons (1.4 onward).
%   \documentclass[11pt]{article}
%   % Shared preamble for the printable mini lessons (1.4 onward).
%   \documentclass[11pt]{article}
%   \input{../preamble}
% Build the key with:  pdflatex -jobname=<name>_KEY "\def\iskey{}\input{<name>}"
% Figures shared with the web page live in figures/Unit_X/<name>.tex and are
% pulled in with \fig{Unit_X/<name>}; the standalone package drops their own
% preamble, so their macros come from figures/preamble.tex, loaded here.
\input{../../figures/preamble}
\usepackage[margin=0.75in]{geometry}
\usepackage{array}
\usepackage{enumitem}
\usepackage{amsmath}
\usepackage[helvet]{sfmath}  % math in the same sans font as the text
\usepackage{standalone}
\pagestyle{empty}
\setlength{\parindent}{0pt}
\setlist[enumerate,1]{leftmargin=1.5em,itemsep=1em}
\setlist[enumerate,2]{label=(\alph*),leftmargin=2em,itemsep=2.2em}

\newif\ifkey
\ifdefined\iskey \keytrue \fi
\colorlet{keyred}{red!75!black}
\newcommand{\ans}[1]{\ifkey\leavevmode{\color{keyred}#1}\fi}
% vertical work space on the student copy, less on the key
\newcommand{\work}[1]{\ifkey\vspace{0.6em}\else\vspace{#1}\fi}
\newcommand{\fig}[1]{\input{../../figures/#1}}
% Title line: \lessontitle{Quartiles, IQR, and Box Plots}
\newcommand{\lessontitle}[1]{{\Large\bfseries #1\ifkey\ \textcolor{keyred}{KEY}\fi}\hfill Name \underline{\hspace{2in}}\par\vspace{0.6em}}
\newcommand{\partner}{\par\vspace{0.4em}Work with a partner. Compare answers when you finish each problem.\par}

% Build the key with:  pdflatex -jobname=<name>_KEY "\def\iskey{}\input{<name>}"
% Figures shared with the web page live in figures/Unit_X/<name>.tex and are
% pulled in with \fig{Unit_X/<name>}; the standalone package drops their own
% preamble, so their macros come from figures/preamble.tex, loaded here.
% Shared preamble for web figures.  Each figure is a standalone .tex that
% \input's this file; scripts/make-figures.sh turns it into an SVG.
\usepackage[T1]{fontenc}
\usepackage{helvet}
\renewcommand{\familydefault}{\sfdefault}
\usepackage{tikz}
\usetikzlibrary{decorations.pathreplacing}
% axis numbers in the text font, not math italic
\pgfkeys{/pgf/number format/.cd,assume math mode=true,1000 sep={}}

% dot plot: \dotplot{min}{max}{value/count,...}{axis label}
\newcommand{\dotplot}[4]{%
\begin{tikzpicture}[x=0.9cm,y=0.33cm]
  \draw[thick] (#1-0.5,0) -- (#2+0.5,0);
  \foreach \t in {#1,...,#2} {\draw (\t,0.25) -- (\t,-0.25) node[below] {\t};}
  \foreach \v/\c in {#3} {\ifnum\c>0 \foreach \k in {1,...,\c} {\fill (\v,\k*1.15) circle (0.16cm);}\fi}
  \node at ({(#1+#2)/2},-2.3) {#4};
\end{tikzpicture}}

% shape sketch histogram (12 bins): \shapehist{heights}{label}
\newcommand{\shapehist}[2]{%
  \foreach \h [count=\i from 0] in {#1} \draw[fill=gray!30] (\i,0) rectangle (\i+1,\h);
  \draw (-0.3,0) -- (12.3,0);
  \node[font=\bfseries] at (6,-2.6) {#2};}
\newcommand{\meanmark}[1]{\fill (#1,0) -- ++(-0.3,-0.75) -- ++(0.6,0) -- cycle;}
\newcommand{\medmark}[2]{\draw[very thick,dashed] (#1,0) -- (#1,#2);}

% number line with minor ticks: \nline{min}{max}{label step}{axis label}
\newcommand{\nline}[4]{%
  \draw[thick] (#1-0.5,0) -- (#2+0.5,0);
  \foreach \t in {#1,...,#2} \draw (\t,0.12) -- (\t,-0.12);
  \pgfmathtruncatemacro{\nxt}{#1+#3}
  \foreach \t in {#1,\nxt,...,#2} \draw (\t,0.2) -- (\t,-0.2) node[below] {\t};
  \node at ({(#1+#2)/2},-0.95) {#4};}
% box plot above the line: \bplot{min}{Q1}{med}{Q3}{max}
\newcommand{\bplot}[5]{%
  \draw[thick] (#1,1.0) -- (#2,1.0);
  \draw[thick] (#4,1.0) -- (#5,1.0);
  \draw[thick] (#1,0.75) -- (#1,1.25);
  \draw[thick] (#5,0.75) -- (#5,1.25);
  \draw[thick,fill=gray!20] (#2,0.55) rectangle (#4,1.45);
  \draw[very thick] (#3,0.55) -- (#3,1.45);}

% ---- added for 1.4 onward ----
% axis on y=0: \axis{min}{max}{label step}{tick step}{axis label}
\newcommand{\axis}[5]{%
  \draw[thick] (#1,0) -- (#2,0);
  \pgfmathsetmacro{\axT}{#1+#4}\foreach \t in {#1,\axT,...,#2} \draw (\t,0.1) -- (\t,-0.1);
  \pgfmathsetmacro{\axL}{#1+#3}\foreach \t in {#1,\axL,...,#2}
    \draw (\t,0.18) -- (\t,-0.18) node[below] {\pgfmathprintnumber{\t}};
  \node at ({(#1+#2)/2},-0.95) {#5};}
% box plot centred at height y: \bplotat{y}{min}{Q1}{med}{Q3}{max}
\newcommand{\bplotat}[6]{%
  \draw[thick] (#2,#1) -- (#3,#1);
  \draw[thick] (#5,#1) -- (#6,#1);
  \draw[thick] (#2,#1-0.22) -- (#2,#1+0.22);
  \draw[thick] (#6,#1-0.22) -- (#6,#1+0.22);
  \draw[thick,fill=gray!20] (#3,#1-0.42) rectangle (#5,#1+0.42);
  \draw[very thick] (#4,#1-0.42) -- (#4,#1+0.42);}
% an outlier dot: \outlier{x}{y}
\newcommand{\outlier}[2]{\fill (#1,#2) circle (2.4pt);}

\usepackage[margin=0.75in]{geometry}
\usepackage{array}
\usepackage{enumitem}
\usepackage{amsmath}
\usepackage[helvet]{sfmath}  % math in the same sans font as the text
\usepackage{standalone}
\pagestyle{empty}
\setlength{\parindent}{0pt}
\setlist[enumerate,1]{leftmargin=1.5em,itemsep=1em}
\setlist[enumerate,2]{label=(\alph*),leftmargin=2em,itemsep=2.2em}

\newif\ifkey
\ifdefined\iskey \keytrue \fi
\colorlet{keyred}{red!75!black}
\newcommand{\ans}[1]{\ifkey\leavevmode{\color{keyred}#1}\fi}
% vertical work space on the student copy, less on the key
\newcommand{\work}[1]{\ifkey\vspace{0.6em}\else\vspace{#1}\fi}
\newcommand{\fig}[1]{\input{../../figures/#1}}
% Title line: \lessontitle{Quartiles, IQR, and Box Plots}
\newcommand{\lessontitle}[1]{{\Large\bfseries #1\ifkey\ \textcolor{keyred}{KEY}\fi}\hfill Name \underline{\hspace{2in}}\par\vspace{0.6em}}
\newcommand{\partner}{\par\vspace{0.4em}Work with a partner. Compare answers when you finish each problem.\par}

% Build the key with:  pdflatex -jobname=<name>_KEY "\def\iskey{}\input{<name>}"
% Figures shared with the web page live in figures/Unit_X/<name>.tex and are
% pulled in with \fig{Unit_X/<name>}; the standalone package drops their own
% preamble, so their macros come from figures/preamble.tex, loaded here.
% Shared preamble for web figures.  Each figure is a standalone .tex that
% \input's this file; scripts/make-figures.sh turns it into an SVG.
\usepackage[T1]{fontenc}
\usepackage{helvet}
\renewcommand{\familydefault}{\sfdefault}
\usepackage{tikz}
\usetikzlibrary{decorations.pathreplacing}
% axis numbers in the text font, not math italic
\pgfkeys{/pgf/number format/.cd,assume math mode=true,1000 sep={}}

% dot plot: \dotplot{min}{max}{value/count,...}{axis label}
\newcommand{\dotplot}[4]{%
\begin{tikzpicture}[x=0.9cm,y=0.33cm]
  \draw[thick] (#1-0.5,0) -- (#2+0.5,0);
  \foreach \t in {#1,...,#2} {\draw (\t,0.25) -- (\t,-0.25) node[below] {\t};}
  \foreach \v/\c in {#3} {\ifnum\c>0 \foreach \k in {1,...,\c} {\fill (\v,\k*1.15) circle (0.16cm);}\fi}
  \node at ({(#1+#2)/2},-2.3) {#4};
\end{tikzpicture}}

% shape sketch histogram (12 bins): \shapehist{heights}{label}
\newcommand{\shapehist}[2]{%
  \foreach \h [count=\i from 0] in {#1} \draw[fill=gray!30] (\i,0) rectangle (\i+1,\h);
  \draw (-0.3,0) -- (12.3,0);
  \node[font=\bfseries] at (6,-2.6) {#2};}
\newcommand{\meanmark}[1]{\fill (#1,0) -- ++(-0.3,-0.75) -- ++(0.6,0) -- cycle;}
\newcommand{\medmark}[2]{\draw[very thick,dashed] (#1,0) -- (#1,#2);}

% number line with minor ticks: \nline{min}{max}{label step}{axis label}
\newcommand{\nline}[4]{%
  \draw[thick] (#1-0.5,0) -- (#2+0.5,0);
  \foreach \t in {#1,...,#2} \draw (\t,0.12) -- (\t,-0.12);
  \pgfmathtruncatemacro{\nxt}{#1+#3}
  \foreach \t in {#1,\nxt,...,#2} \draw (\t,0.2) -- (\t,-0.2) node[below] {\t};
  \node at ({(#1+#2)/2},-0.95) {#4};}
% box plot above the line: \bplot{min}{Q1}{med}{Q3}{max}
\newcommand{\bplot}[5]{%
  \draw[thick] (#1,1.0) -- (#2,1.0);
  \draw[thick] (#4,1.0) -- (#5,1.0);
  \draw[thick] (#1,0.75) -- (#1,1.25);
  \draw[thick] (#5,0.75) -- (#5,1.25);
  \draw[thick,fill=gray!20] (#2,0.55) rectangle (#4,1.45);
  \draw[very thick] (#3,0.55) -- (#3,1.45);}

% ---- added for 1.4 onward ----
% axis on y=0: \axis{min}{max}{label step}{tick step}{axis label}
\newcommand{\axis}[5]{%
  \draw[thick] (#1,0) -- (#2,0);
  \pgfmathsetmacro{\axT}{#1+#4}\foreach \t in {#1,\axT,...,#2} \draw (\t,0.1) -- (\t,-0.1);
  \pgfmathsetmacro{\axL}{#1+#3}\foreach \t in {#1,\axL,...,#2}
    \draw (\t,0.18) -- (\t,-0.18) node[below] {\pgfmathprintnumber{\t}};
  \node at ({(#1+#2)/2},-0.95) {#5};}
% box plot centred at height y: \bplotat{y}{min}{Q1}{med}{Q3}{max}
\newcommand{\bplotat}[6]{%
  \draw[thick] (#2,#1) -- (#3,#1);
  \draw[thick] (#5,#1) -- (#6,#1);
  \draw[thick] (#2,#1-0.22) -- (#2,#1+0.22);
  \draw[thick] (#6,#1-0.22) -- (#6,#1+0.22);
  \draw[thick,fill=gray!20] (#3,#1-0.42) rectangle (#5,#1+0.42);
  \draw[very thick] (#4,#1-0.42) -- (#4,#1+0.42);}
% an outlier dot: \outlier{x}{y}
\newcommand{\outlier}[2]{\fill (#1,#2) circle (2.4pt);}

\usepackage[margin=0.75in]{geometry}
\usepackage{array}
\usepackage{enumitem}
\usepackage{amsmath}
\usepackage[helvet]{sfmath}  % math in the same sans font as the text
\usepackage{standalone}
\pagestyle{empty}
\setlength{\parindent}{0pt}
\setlist[enumerate,1]{leftmargin=1.5em,itemsep=1em}
\setlist[enumerate,2]{label=(\alph*),leftmargin=2em,itemsep=2.2em}

\newif\ifkey
\ifdefined\iskey \keytrue \fi
\colorlet{keyred}{red!75!black}
\newcommand{\ans}[1]{\ifkey\leavevmode{\color{keyred}#1}\fi}
% vertical work space on the student copy, less on the key
\newcommand{\work}[1]{\ifkey\vspace{0.6em}\else\vspace{#1}\fi}
\newcommand{\fig}[1]{\input{../../figures/#1}}
% Title line: \lessontitle{Quartiles, IQR, and Box Plots}
\newcommand{\lessontitle}[1]{{\Large\bfseries #1\ifkey\ \textcolor{keyred}{KEY}\fi}\hfill Name \underline{\hspace{2in}}\par\vspace{0.6em}}
\newcommand{\partner}{\par\vspace{0.4em}Work with a partner. Compare answers when you finish each problem.\par}

% Unit 2 review.  Problems 1-12 follow the Unit 2 quiz question for question
% (same "I ..." goals, two or three times the parts); 13-18 ask the same
% skills in test wording.  2.5 and 2.6 use the printed z-table
% (Z_Table.tex): answers use its 4-decimal areas, and a reverse lookup
% takes the closest area in the table.
\raggedright\hyphenpenalty=10000\exhyphenpenalty=10000
\setlist[enumerate,1]{leftmargin=1.5em,itemsep=1.2em}
\setlist[enumerate,2]{label=(\alph*),leftmargin=2em,itemsep=1.6em}
\newcommand{\lessonhead}[1]{\par\vspace{0.6em}{\large\bfseries #1}\par\vspace{0.25em}\hrule\vspace{0.6em}}
\newcommand{\goal}[1]{\textbf{#1}\par\vspace{0.35em}}
% circled choices: \opt{A} prints plain; \optans{B} is circled in red on the key
\newcommand{\opt}[1]{\tikz[baseline=(c.base)]\node[inner sep=2.5pt](c){#1};}
\newcommand{\optans}[1]{\ifkey\tikz[baseline=(c.base)]\node[circle,inner sep=1.2pt,draw=keyred,thick,text=keyred](c){#1};\else\opt{#1}\fi}
\newcommand{\choices}[4]{\par\vspace{0.3em}\begin{tabular}{@{}l@{\hspace{0.6em}}p{6in}}#1\\#2\\#3\\#4\\\end{tabular}}
% a circled pair of measures (oval on the key): \pair{mean and SD} / \pairans{median and IQR}
\newcommand{\pair}[1]{\tikz[baseline=(c.base)]\node[rounded corners=0.6em,inner sep=2.5pt](c){#1};}
\newcommand{\pairans}[1]{\ifkey\tikz[baseline=(c.base)]\node[rounded corners=0.6em,inner sep=2.5pt,draw=keyred,thick,text=keyred](c){#1};\else\pair{#1}\fi}
% one row of an SD table (as in 2.1): \sdrow{x}{x - mean}{(x - mean)^2}; the last two show on the key only
\newcommand{\sdrow}[3]{#1 & \ans{#2} & \ans{#3} \\ \hline}
\newenvironment{sdtable}{\renewcommand{\arraystretch}{1.45}\begin{tabular}{|c|p{1.0in}|p{1.0in}|}\hline
  $x$ & $x - \bar{x}$ & $(x - \bar{x})^2$ \\ \hline}{\end{tabular}}
% blank normal curve over a number line; the key writes the values in red:
% \labelcurve{v1/v2/.../v7 as z/value pairs}{axis label}
\newcommand{\labelcurve}[2]{%
\begin{tikzpicture}[x=1.3cm,y=0.7cm]
  \bellcurve
  \draw[thick] (-3.8,0) -- (3.8,0);
  \foreach \z in {-3,...,3} \draw (\z,0.12) -- (\z,-0.12);
  \ifkey\begin{scope}[keyred]\foreach \z/\v in {#1} \node[below] at (\z,-0.12) {\v};\end{scope}\fi
  \node at (0,-1.3) {#2};
\end{tikzpicture}}

\begin{document}

\lessontitle{Unit 2 Review}

\textbf{One study, every tool.} A class flew paper airplanes. Here are 5 flights, in feet.
\begin{center}14,\; 18,\; 22,\; 22,\; 24\end{center}

\vspace{-0.4em}
\renewcommand{\arraystretch}{1.3}
\begin{tabular}{|>{\raggedright\arraybackslash}p{1.8in}|>{\raggedright\arraybackslash}p{4.7in}|}
\hline
\textbf{2.1 Find the SD.}

The SD is about how far the values usually are from the mean. &
$\bar{x} = 100 \div 5 = 20$ feet

\vspace{0.2em}
\centerline{\renewcommand{\arraystretch}{1.15}\small
\begin{tabular}{|c|c|c|c|c|c|}\hline
$x$ & 14 & 18 & 22 & 22 & 24 \\ \hline
$x - \bar{x}$ & $-6$ & $-2$ & 2 & 2 & 4 \\ \hline
$(x - \bar{x})^2$ & 36 & 4 & 4 & 4 & 16 \\ \hline
\end{tabular}}

\vspace{0.3em}
$\Sigma(x - \bar{x})^2 = 64$ \qquad $s = \sqrt{\dfrac{\Sigma(x - \bar{x})^2}{n - 1}} = \sqrt{\dfrac{64}{4}} = \sqrt{16} = 4$ feet

The flights are usually about 4 feet from the mean. \\ \hline
\textbf{2.2 Choose the measures.}

Symmetric, no outliers: mean and SD.\newline Skewed or outliers: median and IQR. &
A sixth plane flew 62 feet.
\centerline{\fig{Unit_2/Review_planes_dot_plot}}
Mean $= 162 \div 6 = 27$ feet. \quad SD $\approx 17.5$ feet.\newline
Median $= 22$ feet. \quad Q1 $= 18$, \; Q3 $= 24$, \; IQR $= 24 - 18 = 6$ feet.

Upper fence $= 24 + 1.5 \times 6 = 33$, so 62 is an outlier. It pulled the mean from 20 up to 27. Use the \textbf{median and the IQR}. \\ \hline
\textbf{2.3 Where does one value stand?}

Percentile: the percent of the data below the value.\newline
z-score: how many SDs the value is above or below the mean. &
4 of the 5 flights are shorter than 24 feet. $4 \div 5 = 0.80$. So 24 feet is at the \textbf{80th percentile}.

All the school's flights: mean $\mu = 20$ feet, SD $\sigma = 4$ feet.\newline
A 26-foot flight: \; $z = \dfrac{x - \mu}{\sigma} = \dfrac{26 - 20}{4} = 1.5$. \; It is 1.5 SDs above the mean. \\ \hline
\textbf{2.4 Use the empirical rule.}

In a normal distribution, about 68\% is within 1 SD of the mean, 95\% within 2 SDs, and 99.7\% within 3 SDs. &
The school's flights have a normal distribution.
\centerline{\fig{Unit_2/Review_planes_normal}}
68\% of the flights are from 16 to 24 feet. 95\% are from 12 to 28 feet. \\ \hline
\textbf{2.5 Use the z-table.}

The table gives the area below $z$. &
\rule{0pt}{2.4em}Shorter than 25 feet: \; $z = \dfrac{25 - 20}{4} = 1.25$. \; Table: .8944. \; \textbf{89.44\%}\newline
Longer than 25 feet: \; $1 - .8944 = .1056$. \; \textbf{10.56\%} \\ \hline
\textbf{2.6 Work in reverse.}

Find the area inside the table. Read $z$. Then $x = \mu + z\sigma$. &
The longest 10\% of flights win a prize. Area to the left: $1 - .10 = .90$.\newline
Closest area .8997, at $z = 1.28$. \; $x = 20 + 1.28 \times 4 = 25.12$. \; A flight must go about \textbf{25.1 feet}. \\ \hline
\end{tabular}

\vspace{0.5em}
\textbf{Above or top? Subtract from 1.} The z-table always gives the area below $z$ (to the left).

\partner

% ---------------- PROBLEMS ----------------
\newpage
\lessonhead{2.1 - Standard Deviation}
\begin{enumerate}

\item \goal{I find the standard deviation of a small data set.}
A bike shop sold this many bikes each day for 5 days.
\begin{center}4,\; 11,\; 7,\; 11,\; 7\end{center}
\begin{minipage}[t]{0.5\linewidth}\vspace{0pt}\raggedright
\begin{enumerate}
  \item Find the mean, $\bar{x}$. Then complete the table.

  \vspace{0.5em}
  $\bar{x} =$ \underline{\makebox[1.3in]{\ans{$40 \div 5 = 8$}}}

  \vspace{0.6em}
  \begin{sdtable}
  \sdrow{4}{$-4$}{16}
  \sdrow{11}{3}{9}
  \sdrow{7}{$-1$}{1}
  \sdrow{11}{3}{9}
  \sdrow{7}{$-1$}{1}
  \multicolumn{2}{|r|}{$\Sigma(x - \bar{x})^2$} & \ans{36} \\ \hline
  \end{sdtable}
\end{enumerate}
\end{minipage}\hfill
\begin{minipage}[t]{0.47\linewidth}\vspace{0pt}\raggedright
\begin{enumerate}[start=2]
  \item Find the standard deviation, $s$. Round to 2 decimal places if you need to.

  \vspace{0.6em}
  \ans{$s = \sqrt{\frac{\Sigma(x - \bar{x})^2}{n - 1}} = \sqrt{\frac{36}{4}} = \sqrt{9} = 3$ bikes}
\end{enumerate}
\end{minipage}

\vspace{1.4em}
Four friends scored these points on a 10-point quiz.
\begin{center}9,\; 5,\; 8,\; 6\end{center}
\begin{minipage}[t]{0.5\linewidth}\vspace{0pt}\raggedright
\begin{enumerate}[start=3]
  \item Find the mean, $\bar{x}$. Then complete the table.

  \vspace{0.5em}
  $\bar{x} =$ \underline{\makebox[1.3in]{\ans{$28 \div 4 = 7$}}}

  \vspace{0.6em}
  \begin{sdtable}
  \sdrow{9}{2}{4}
  \sdrow{5}{$-2$}{4}
  \sdrow{8}{1}{1}
  \sdrow{6}{$-1$}{1}
  \multicolumn{2}{|r|}{$\Sigma(x - \bar{x})^2$} & \ans{10} \\ \hline
  \end{sdtable}
\end{enumerate}
\end{minipage}\hfill
\begin{minipage}[t]{0.47\linewidth}\vspace{0pt}\raggedright
\begin{enumerate}[start=4]
  \item Find the standard deviation, $s$. Round to 2 decimal places.

  \vspace{0.6em}
  \ans{$s = \sqrt{\frac{10}{4 - 1}} = \sqrt{\frac{10}{3}} \approx \sqrt{3.33} \approx 1.83$ points}
\end{enumerate}
\end{minipage}

\newpage
\item \goal{I compare the spread of data sets by looking at graphs, and I explain what the standard deviation tells about data.}
Each dot plot has 10 values and a mean of 6.
\begin{center}\fig{Unit_2/Review_sd_three_plots}\end{center}
\begin{enumerate}[itemsep=1.2em]
  \item Put the dot plots in order from the \textbf{smallest} SD to the \textbf{largest} SD. \ans{\quad L, J, K}
  \item Explain your order.

  \ans{The values in L are packed close to 6. The values in K are far from 6.\\
  (The SDs are about 0.82, 1.83, and 3.94.)}
  \work{2.5em}
  \item The SD of plot K is about 3.94. What does that number tell you?

  \ans{The values in K are usually about 4 away from the mean of 6.}
  \work{2.5em}
  \item A data set has an SD of 0. What do you know about its values?

  \ans{All the values are the same. Every value is equal to the mean.}
  \work{2em}
  \item Set M is 40, 45, 50, 55, 60. Set N is 48, 49, 50, 51, 52. Both have a mean of 50. Without calculating, which set has the larger SD? Explain.

  \ans{Set M. Its values are farther from 50.}
  \work{2em}
  \item Can a standard deviation be negative? Explain.

  \ans{No. It comes from adding squares (never negative) and taking a square root.}
\end{enumerate}

\end{enumerate}

% ======================================================================
\newpage
\lessonhead{2.2 - Choosing Measures and Comparing Distributions}
\begin{enumerate}[resume]

\item \goal{I choose mean and SD or median and IQR based on the shape of the data.}
For each data set, circle the measures you should use. Then explain.
\begin{enumerate}[itemsep=1.1em]
  \item The heights of 500 corn plants in a field. The histogram is symmetric, with no outliers.

  \pairans{mean and SD}\hspace{1em}\pair{median and IQR}\hspace{1.5em}\ans{The data is symmetric with no outliers.}
  \work{2em}
  \item The number of followers 40 students have on a social media app. The histogram is skewed right.

  \pair{mean and SD}\hspace{1em}\pairans{median and IQR}\hspace{1.5em}\ans{The data is skewed.}
  \work{2em}
  \item Scores on a very easy test. The histogram is skewed left.

  \pair{mean and SD}\hspace{1em}\pairans{median and IQR}\hspace{1.5em}\ans{The data is skewed.}
  \work{2em}
  \item The times 200 students took to run a mile. The histogram is symmetric, with no outliers.

  \pairans{mean and SD}\hspace{1em}\pair{median and IQR}\hspace{1.5em}\ans{The data is symmetric with no outliers.}
  \work{2em}
  \item Minutes 7 students spent on homework: \; 30,\; 35,\; 40,\; 40,\; 45,\; 50,\; 120

  \pair{mean and SD}\hspace{1em}\pairans{median and IQR}\hspace{1.5em}\ans{120 is an outlier.}

  \ans{Q1 $= 35$, Q3 $= 50$, IQR $= 15$. Upper fence $= 50 + 1.5 \times 15 = 72.5$, and 120 is above it.}
  \work{2.5em}
  \item High temperatures for 9 days, in $^\circ$F: \; 70,\; 72,\; 73,\; 74,\; 75,\; 76,\; 77,\; 78,\; 80

  \pairans{mean and SD}\hspace{1em}\pair{median and IQR}\hspace{1.5em}\ans{Symmetric (mean $=$ median $= 75$), no outliers.}

  \ans{Q1 $= 72.5$, Q3 $= 77.5$, IQR $= 5$. Fences: 65 and 85. No value is outside them.}
  \work{2.5em}
  \item For the homework data in (e), find the mean and the median. Which one better describes a typical student? Explain.

  \ans{Mean $= 360 \div 7 \approx 51.4$ minutes. Median $= 40$ minutes. The median is better. 120 pulls the mean up, and 6 of the 7 students spent less than the mean.}
  \work{3.5em}
\end{enumerate}

\newpage
\item \goal{I compare two distributions using shape, center, spread, and outliers.}
These dot plots show how many hours students in two clubs volunteered last month. Each club has 20 students.

\vspace{0.3em}
\begin{minipage}[t]{0.57\linewidth}\vspace{0pt}\raggedright
\fig{Unit_2/Review_volunteer_dot_plots}
\end{minipage}\hfill
\begin{minipage}[t]{0.4\linewidth}\vspace{0.8em}\raggedright
\renewcommand{\arraystretch}{1.3}
\begin{tabular}{|l|c|c|c|c|}\hline
 & Q1 & Median & Q3 & IQR \\ \hline
Art & 4 & 5.5 & 7 & 3 \\ \hline
Science & 1 & 2 & 3.5 & 2.5 \\ \hline
\end{tabular}
\end{minipage}

\vspace{0.4em}
\begin{enumerate}[itemsep=0.7em]
  \item Describe the shape of each club's data. \ans{\quad Art: symmetric. \; Science: skewed right.}
  \work{1em}
  \item Compare the centers. Use numbers and the context.

  \ans{The median for the art club (5.5 hours) is greater than for the science club (2 hours). Art club students usually volunteered 3.5 hours more.}
  \work{2.5em}
  \item Compare the spreads. Use numbers.

  \ans{The IQR for the art club (3 hours) is a little greater than for the science club (2.5 hours). The middle half of the art club is a little more spread out.}
  \work{2.5em}
  \item Does either club have an outlier? Explain.

  \ans{Yes. One science club student volunteered 13 hours. The upper fence is $3.5 + 1.5 \times 2.5 = 7.25$, and 13 is above it.}
  \work{2.5em}
  \item Should you compare these clubs with the mean and SD or with the median and IQR? Why?

  \ans{The median and IQR. The science club is skewed right and has an outlier.}
\end{enumerate}

\vspace{0.8em}
Two pizza places recorded their delivery times. Both histograms are symmetric.
\begin{center}
\renewcommand{\arraystretch}{1.3}
\begin{tabular}{|l|c|c|}\hline
 & Mean & SD \\ \hline
Place A & 30 minutes & 3 minutes \\ \hline
Place B & 26 minutes & 8 minutes \\ \hline
\end{tabular}
\end{center}
\begin{enumerate}[start=6,itemsep=0.7em]
  \item Which place is faster on average? \ans{\quad Place B, by 4 minutes.}
  \item Which place has delivery times that are more spread out? \ans{\quad Place B. Its SD is larger.}
  \item Write one sentence that compares the two places. Use both the center and the spread.

  \ans{Place B is faster on average (26 vs 30 minutes), but its times are more spread out (SD 8 vs 3 minutes).}
\end{enumerate}

\end{enumerate}

% ======================================================================
\newpage
\lessonhead{2.3 - Percentiles and Z-Scores}
\begin{enumerate}[resume]

\item \goal{I find the percentile of a value.}
20 students said how many minutes they practiced a musical instrument yesterday. Here are the times, in order.
\begin{center}10,\; 12,\; 15,\; 18,\; 20,\; 24,\; 25,\; 28,\; 30,\; 32,\; 35,\; 36,\; 40,\; 42,\; 45,\; 48,\; 50,\; 55,\; 60,\; 90\end{center}
\begin{enumerate}[itemsep=1.2em]
  \item What percentile is 40 minutes? \ans{\quad 12 times are below 40. \; $12 \div 20 = 0.60$, the 60th percentile}
  \work{1em}
  \item What percentile is 15 minutes? \ans{\quad 2 times are below 15. \; $2 \div 20 = 0.10$, the 10th percentile}
  \work{1em}
  \item Which time is at the 25th percentile? \ans{\quad 25\% of 20 is 5 times below it. The 6th time: \textbf{24 minutes}}
  \work{1em}
  \item Which time is at the 90th percentile? \ans{\quad 90\% of 20 is 18 times below it. The 19th time: \textbf{60 minutes}}
  \work{1em}
\end{enumerate}

\vspace{0.6em}
This dot plot shows how many goals 25 hockey players scored this season.
\begin{center}\fig{Unit_2/Review_goals_dot_plot}\end{center}
\begin{enumerate}[start=5,itemsep=1.2em]
  \item What percentile is 6 goals? \ans{\quad $2 + 3 + 4 + 5 + 4 = 18$ below. \; $18 \div 25 = 0.72$, the 72nd percentile}
  \work{1em}
  \item What percentile is 3 goals? \ans{\quad $2 + 3 = 5$ below. \; $5 \div 25 = 0.20$, the 20th percentile}
  \work{1em}
  \item Nadia says, ``My test score is at the 90th percentile, so I got 90\% of the questions right.'' Explain her mistake.

  \ans{A percentile is not a percent score. She scored higher than 90\% of the students. It does not tell what percent of the questions she got right.}
\end{enumerate}

\newpage
\item \goal{I find and explain a z-score, and I compare values from different data sets with z-scores.}
Scores on a spelling test have a mean $\mu = 50$ and an SD $\sigma = 6$. Find the z-score of each score. Then say what it means.
\begin{enumerate}[itemsep=1.2em]
  \item 62 \ans{\quad $z = \frac{62 - 50}{6} = 2$. \; 2 SDs above the mean.}
  \work{1.5em}
  \item 47 \ans{\quad $z = \frac{47 - 50}{6} = -0.5$. \; Half an SD below the mean.}
  \work{1.5em}
  \item 50 \ans{\quad $z = \frac{50 - 50}{6} = 0$. \; Equal to the mean.}
  \work{1.5em}
  \item 41 \ans{\quad $z = \frac{41 - 50}{6} = -1.5$. \; 1.5 SDs below the mean.}
  \work{1.5em}
\end{enumerate}

\vspace{0.6em}
Omar and Lena compete in different events.
\begin{center}
\renewcommand{\arraystretch}{1.3}
\begin{tabular}{|l|l|c|c|c|}\hline
Athlete & Event & Jump & Mean $\mu$ & SD $\sigma$ \\ \hline
Omar & long jump & 5.6 m & 5.0 m & 0.3 m \\ \hline
Lena & high jump & 1.55 m & 1.40 m & 0.10 m \\ \hline
\end{tabular}
\end{center}
\begin{enumerate}[start=5,itemsep=1.2em]
  \item Find each athlete's z-score.

  \ans{Omar: $z = \frac{5.6 - 5.0}{0.3} = \frac{0.6}{0.3} = 2$ \qquad Lena: $z = \frac{1.55 - 1.40}{0.10} = \frac{0.15}{0.10} = 1.5$}
  \work{3em}
  \item Lena jumped 1.55 m and Omar jumped 5.6 m. Who did better compared to the others in the event? Explain.

  \ans{Omar. His z-score (2) is larger than Lena's (1.5). Their jumps are in different events, so we compare z-scores, not meters.}
  \work{2.5em}
\end{enumerate}

\vspace{0.6em}
Go back to the spelling test ($\mu = 50$, $\sigma = 6$). Use $x = \mu + z\sigma$.
\begin{enumerate}[start=7,itemsep=1.2em]
  \item What score has a z-score of 1.5? \ans{\quad $50 + 1.5 \times 6 = 50 + 9 = 59$}
  \work{1.5em}
  \item What score has a z-score of $-2$? \ans{\quad $50 + (-2)(6) = 50 - 12 = 38$}
\end{enumerate}

\end{enumerate}

% ======================================================================
\newpage
\lessonhead{2.4 - Normal Distributions and the Empirical Rule}
\begin{enumerate}[resume]

\item \goal{I use the empirical rule (68-95-99.7) to find percents, and I find the percentile of a value 1, 2, or 3 SDs from the mean.}
The heights of sunflowers in a field have a normal distribution. The mean is 180 cm and the SD is 12 cm.
\begin{enumerate}[itemsep=1.1em]
  \item Label the number line with the heights.

  \begin{center}\labelcurve{-3/144,-2/156,-1/168,0/180,1/192,2/204,3/216}{Height of a sunflower (cm)}\end{center}
  \item About what percent of the sunflowers are from 168 to 192 cm tall? \ans{\quad 68\%}
  \item About what percent are from 156 to 204 cm tall? \ans{\quad 95\%}
  \item About what percent are taller than 192 cm? \ans{\quad $13.5 + 2.35 + 0.15 = 16\%$}
  \item About what percent are shorter than 156 cm? \ans{\quad $2.35 + 0.15 = 2.5\%$}
  \item About what percent are from 180 to 204 cm tall? \ans{\quad $34 + 13.5 = 47.5\%$}
  \item About what percent are from 192 to 216 cm tall? \ans{\quad $13.5 + 2.35 = 15.85\%$}
  \item What percentile is a height of 204 cm? \ans{\quad 2 SDs above. \; $100 - 2.5 = 97.5$, the 97.5th percentile}
  \item What percentile is a height of 168 cm? \ans{\quad 1 SD below. \; $0.15 + 2.35 + 13.5 = 16$, the 16th percentile}
\end{enumerate}

\newpage
\item \goal{I describe a normal distribution, and I compare normal curves with different means and SDs.}
The three curves are normal distributions.
\begin{center}\fig{Unit_2/Review_three_curves}\end{center}
\begin{enumerate}[itemsep=1.1em]
  \item Which two curves have the same mean? \ans{\quad X and Y (both 12)}
  \item Which two curves have the same SD? \ans{\quad X and Z (same width and shape)}
  \item Which curve has the smallest SD? Explain. \ans{\quad Y. It is the tallest and the narrowest.}
  \work{1em}
  \item How are X and Z the same? How are they different?

  \ans{Same SD (same width and shape). Different means: X is centered at 12, Z at 24.}
  \work{2em}
\end{enumerate}

\vspace{0.4em}
Would each data set have a normal distribution? Explain.
\begin{enumerate}[start=5,itemsep=1.1em]
  \item The heights of 1,000 adult women. \ans{\quad Yes. Heights are bell-shaped and symmetric.}
  \work{1em}
  \item The prices of every car sold in a city. \ans{\quad No. A few very expensive cars make it skewed right.}
  \work{1em}
  \item Scores on a very hard test. \ans{\quad No. Most scores are low, so it is skewed right.}
  \work{1em}
  \item Which statement is true for every normal distribution? Circle the letter.
  \choices{\opt{A} & About 95\% of the data is within 1 SD of the mean.}
          {\opt{B} & A larger SD makes the curve taller and narrower.}
          {\optans{C} & The curve is symmetric about the mean.}
          {\opt{D} & The median is greater than the mean.}
\end{enumerate}

\end{enumerate}

% ======================================================================
\newpage
\lessonhead{2.5 - Normal Calculations with a Z-Table}
Use the z-table.
\begin{enumerate}[resume]

\item \goal{I read a z-table.}
\begin{minipage}[t]{0.48\linewidth}\vspace{0pt}\raggedright
Find the area \textbf{below} each z-score.
\begin{enumerate}[itemsep=1em]
  \item $z = 0.36$ \ans{\quad .6406}
  \item $z = -2.15$ \ans{\quad .0158}
  \item $z = 1.82$ \ans{\quad .9656}
  \item $z = -0.07$ \ans{\quad .4721}
\end{enumerate}
\end{minipage}\hfill
\begin{minipage}[t]{0.48\linewidth}\vspace{0pt}\raggedright
Find the area \textbf{above} each z-score.
\begin{enumerate}[start=5,itemsep=1em]
  \item $z = 0.95$ \ans{\quad $1 - .8289 = .1711$}
  \item $z = -1.38$ \ans{\quad $1 - .0838 = .9162$}
\end{enumerate}

\vspace{1.2em}
Find the area \textbf{between} the two z-scores.
\begin{enumerate}[start=7,itemsep=1em]
  \item $z = -2$ and $z = 2$ \ans{\quad $.9772 - .0228 = .9544$}
  \item $z = 0.25$ and $z = 1.75$ \ans{\quad $.9599 - .5987 = .3612$}
\end{enumerate}
\end{minipage}

\vspace{0.6em}
\item \goal{I find the percent of data below, above, or between values in a normal distribution.}
The lengths of movies at a theater have a normal distribution with a mean of 110 minutes and an SD of 12 minutes. What percent of the movies are:
\begin{enumerate}[itemsep=0.8em]
  \item shorter than 125 minutes?

  \ans{$z = \frac{125 - 110}{12} = 1.25$. \; Table: .8944. \; \textbf{89.44\%}}
  \work{3.5em}
  \item longer than 128 minutes?

  \ans{$z = \frac{128 - 110}{12} = 1.5$. \; $1 - .9332 = .0668$. \; \textbf{6.68\%}}
  \work{3.5em}
  \item from 100 to 120 minutes long?

  \ans{$z = \frac{100 - 110}{12} = -0.83$, table .2033. \; $z = \frac{120 - 110}{12} = 0.83$, table .7967.\\[0.5em] $.7967 - .2033 = .5934$. \; \textbf{59.34\%}}
  \work{3.5em}
\end{enumerate}

\ifkey\newpage\else\vspace{0.6em}\fi
Bags of trail mix have weights with a normal distribution. The mean is 250 grams and the SD is 4 grams. What percent of the bags weigh:
\begin{enumerate}[start=4,itemsep=0.8em]
  \item less than 245 grams?

  \ans{$z = \frac{245 - 250}{4} = -1.25$. \; Table: .1056. \; \textbf{10.56\%}}
  \work{3.5em}
  \item more than 251 grams?

  \ans{$z = \frac{251 - 250}{4} = 0.25$. \; $1 - .5987 = .4013$. \; \textbf{40.13\%}}
  \work{3.5em}
  \item from 247 to 256 grams?

  \ans{$z = \frac{247 - 250}{4} = -0.75$, table .2266. \; $z = \frac{256 - 250}{4} = 1.5$, table .9332.\\[0.5em] $.9332 - .2266 = .7066$. \; \textbf{70.66\%}}
  \work{3.5em}
\end{enumerate}

\end{enumerate}

% ======================================================================
\newpage
\lessonhead{2.6 - Normal Calculations in Reverse}
Use the z-table.
\begin{enumerate}[resume]

\item \goal{I find the z-score for a percentile with the z-table, and I find the value at a percentile in a normal distribution.}
Find the z-score at each percentile.

\begin{minipage}[t]{0.48\linewidth}\vspace{0pt}\raggedright
\begin{enumerate}[itemsep=1.2em]
  \item 72nd percentile \ans{\newline closest area .7190: \; $z = 0.58$}
  \item 15th percentile \ans{\newline closest area .1492: \; $z = -1.04$}
\end{enumerate}
\end{minipage}\hfill
\begin{minipage}[t]{0.48\linewidth}\vspace{0pt}\raggedright
\begin{enumerate}[start=3,itemsep=1.2em]
  \item 99th percentile \ans{\newline closest area .9901: \; $z = 2.33$}
  \item 42nd percentile \ans{\newline closest area .4207: \; $z = -0.20$}
\end{enumerate}
\end{minipage}

\vspace{0.8em}
The times to finish one level of a puzzle game have a normal distribution. The mean is 90 seconds and the SD is 15 seconds.
\begin{enumerate}[start=5,itemsep=1em]
  \item What time is at the 72nd percentile?

  \ans{$z = 0.58$. \; $x = \mu + z\sigma = 90 + 0.58 \times 15 = 90 + 8.7 = 98.7$ seconds}
  \work{3em}
  \item What time is at the 15th percentile?

  \ans{$z = -1.04$. \; $90 + (-1.04)(15) = 90 - 15.6 = 74.4$ seconds}
  \work{3em}
\end{enumerate}

\item \goal{I find cutoffs for ``top'' and ``bottom'' percents.}
Typing speeds in a class have a normal distribution. The mean is 40 words per minute and the SD is 6 words per minute.
\begin{enumerate}[itemsep=1em]
  \item The top 2\% get a certificate. How fast must a student type to get one?

  \ans{Top 2\%: the area to the left is $1 - .02 = .98$. Closest area .9798: \; $z = 2.05$.\\
  $40 + 2.05 \times 6 = 40 + 12.3 = 52.3$ words per minute}
  \work{3em}
  \item Students in the bottom 25\% get extra practice. Below what speed does a student get extra practice?

  \ans{Bottom 25\%: the area is .25. Closest area .2514: \; $z = -0.67$.\\
  $40 + (-0.67)(6) = 40 - 4.02 = 35.98$. About \textbf{36.0 words per minute}}
  \work{3em}
\end{enumerate}

\vspace{0.6em}
The weights of pumpkins at a farm have a normal distribution. The mean is 12 pounds and the SD is 2 pounds.
\begin{enumerate}[start=3,itemsep=1em]
  \item The heaviest 10\% are sold as ``giant'' pumpkins. How heavy must a pumpkin be to be giant?

  \ans{Top 10\%: the area to the left is $1 - .10 = .90$. Closest area .8997: \; $z = 1.28$.\\
  $12 + 1.28 \times 2 = 12 + 2.56 = 14.56$ pounds}
  \work{3em}
  \item The lightest 20\% are used to make pies. Below what weight is a pumpkin used for pie?

  \ans{Bottom 20\%: the area is .20. Closest area .2005: \; $z = -0.84$.\\
  $12 + (-0.84)(2) = 12 - 1.68 = 10.32$ pounds}
  \work{3em}
\end{enumerate}

\end{enumerate}

% ======================================================================
\newpage
\lessonhead{Test-Style Questions}
A test asks the same questions with more words. Read slowly. Circle the best answer.

\begin{enumerate}[resume,itemsep=1.2em]

\item Each data set below has a mean of 50. Which data set has the greatest standard deviation?
\choices{\opt{A} & 48, 49, 50, 51, 52}
        {\optans{B} & 30, 40, 50, 60, 70}
        {\opt{C} & 45, 50, 50, 50, 55}
        {\opt{D} & 50, 50, 50, 50, 50}

\item A restaurant recorded the number of minutes that each of 50 customers waited for a table. The distribution of the wait times is strongly skewed to the right and has several outliers. Which measures of center and spread are most appropriate to describe the wait times?
\choices{\optans{A} & the median and the interquartile range}
        {\opt{B} & the mean and the standard deviation}
        {\opt{C} & the mean and the interquartile range}
        {\opt{D} & the median and the standard deviation}

\item On a science test, the class mean was 72 points with a standard deviation of 6 points. On a history test, the class mean was 80 points with a standard deviation of 5 points. Mia scored 84 points on the science test and 88 points on the history test. Which statement is true?
\choices{\opt{A} & Mia did better on the history test compared to her class, because 88 is greater than 84.}
        {\opt{B} & Mia did equally well on both tests compared to her class, because both scores are above the class mean.}
        {\opt{C} & Mia's history score was 1.6 points above the class mean.}
        {\optans{D} & Mia did better on the science test compared to her class, because her z-score was greater on that test.}
\par\ans{Science: $z = \frac{84 - 72}{6} = 2$. \; History: $z = \frac{88 - 80}{5} = 1.6$. \; C is wrong: 88 is 8 points (1.6 SDs) above the mean.}

\newpage

\item The lifetimes of a type of light bulb are approximately normally distributed with a mean of 900 hours and a standard deviation of 50 hours. According to the empirical rule, approximately what percent of the light bulbs last between 800 hours and 950 hours?
\choices{\opt{A} & 47.5\%}
        {\opt{B} & 68\%}
        {\optans{C} & 81.5\%}
        {\opt{D} & 95\%}
\par\ans{800 is 2 SDs below the mean and 950 is 1 SD above. \; $13.5 + 34 + 34 = 81.5\%$}

\item The scores on a math contest are approximately normally distributed with a mean of 60 points and a standard deviation of 8 points. Use the z-table. What proportion of the scores are greater than 74 points?
\choices{\optans{A} & .0401}
        {\opt{B} & .4599}
        {\opt{C} & .5401}
        {\opt{D} & .9599}
\par\ans{$z = \frac{74 - 60}{8} = 1.75$. \; $1 - .9599 = .0401$}

\item The weights of watermelons at a farm are approximately normally distributed with a mean of 20 pounds and a standard deviation of 2.5 pounds. The heaviest 3\% of the watermelons win a blue ribbon. Use the z-table. Which weight is closest to the lightest watermelon that wins a blue ribbon?
\choices{\opt{A} & 15.3 pounds}
        {\opt{B} & 22.5 pounds}
        {\optans{C} & 24.7 pounds}
        {\opt{D} & 25.0 pounds}
\par\ans{Area to the left $= 1 - .03 = .97$. Closest area .9699: \; $z = 1.88$. \; $20 + 1.88 \times 2.5 = 24.7$}

\end{enumerate}

\end{document}
